% GDZ PPN236018647,native565lower–567;HKUST566–568.
\NumberedItem{95}
\indent Novus in analysi campus se nobis aperuit, scilicet investigatio functionum etc.
\par\noindent\hfill[1798] Oct.

\medskip\begingroup\small
\indent Nach den bei der Nr. 94 erwähnten Aufzeichnungen finden wir auf S. 21 ff. der Scheda Aa wieder Entwicklungen zur Theorie der lemniskatischen Funktionen, die man mit den hier angekündigten Fortschritten in Verbindung zu setzen hat. Was auf S. 21, 22 der Scheda steht, ist Werke III, S. 418, soweit noch art. [7.] reicht, abgedruckt. Es folgen dann S. 23—25 der Scheda Rechnungen, die darauf hinzielen, aus den Entwicklungen von \(\log P,\ \log Q\) analoge trigonometrische Reihenentwicklungen für \(P,Q\) selbst zu
\par\endgroup

\clearpage\SourceLeaf{566}\ArithmeticHead{537}{R}{1798 IUL.—1798 OCT.}
\begingroup\small
\makeatletter\tagsleft@true\let\veqno\@@leqno\makeatother
\noindent erhalten. Es gelingt dies zunächst in der Weise, daß die Koeffizienten numerisch gefunden werden; so steht z. B. auf S. 25
\[
\begin{array}{r@{\;}r@{\quad}r@{}c@{}l@{\;}l}
P\varphi={}&0{,}8346268416&7407316&.&4&\sin\varphi\\
{}+{}&62694861&2274007&.&5&(\sin\varphi)^3\\
{}+{}&874&5689900&.&5&(\sin\varphi)^5\\
{}+{}&&238&.&56&(\sin\varphi)^7,
\end{array}
\]
\noindent wo, wie S. 24 bemerkt wird, der Koeffizient von \(\sin\varphi\) nichts anderes ist als \(\frac\varpi\pi\), und daraus wird nun weiter die Form
\[
\begin{array}{r@{\;}r@{\quad}r@{\;}l}
P\varphi={}&0{,}8393290109&26691403&\sin\varphi\\
{}-{}&15673988&60966741&\sin3\varphi\\
{}+{}&54&66056449&\sin5\varphi\\
{}-{}&&\ldots37&\sin7\varphi
\end{array}
\]
\noindent gewonnen. Das Gesetz der Koeffizienten in dieser Entwicklung und in der analogen für \(Q\varphi\) scheint \textsc{Gauss} durch numerische Induktion gefunden zu haben, das zeigen die auf S. 29—33 befindlichen Rechnungen für \(e^{-\frac14\pi},\ e^{-\pi},\ e^{-\frac94\pi}\). Auf S. 27 und 28 stehen die Werke III, S. 418 im art. [8.] abgedruckten Formeln für \(\operatorname{sin\,lem}\varphi,\ P\varphi,\ Q\varphi\)\footnote{Auch bei diesen ist im Abdruck beiderseits auf Bogenmaß reduziert.}. Diese endgültigen Darstellungen von \(P,Q\):
\begin{equation}
\left\{\begin{aligned}
P\varphi&=2^{\frac34}\surd\frac\pi\varpi\left(\frac{\sin\varphi}{e^{\frac14\pi}}-\frac{\sin3\varphi}{e^{\frac94\pi}}+\frac{\sin5\varphi}{e^{\frac{25}4\pi}}-\ldots\right),\\[1ex]
Q\varphi&=\frac{\sqrt{\frac\pi\varpi}}{2^{\frac14}}\left(1+\frac{2\cos2\varphi}{e^\pi}+\frac{2\cos4\varphi}{e^{4\pi}}+\ldots\right)
\end{aligned}\right.\tag*{(*)}
\end{equation}
\noindent hat \textsc{Gauss} wohl durch Einsetzen in die auf S. 23 der Scheda aufgezeichnete Gleichung
\[
\begin{aligned}
\frac d{[d\varphi]}\frac PQ&=\frac1Q\frac{dP}{d\varphi}-\frac P{QQ}\frac{dQ}{d\varphi}=\surd\left(1-\frac{P^4}{Q^4}\right)\\
QP'-PQ'&=\surd(Q^4-P^4)
\end{aligned}
\]
\noindent bestätigt (vergl. auch die Leistenotiz oben S. 167, art. [2.]).

Die angegebenen Reihen für \(P,Q\) liefern auch sofort die S. 28 der Scheda unmittelbar anschließend (wie in dem Abdruck Werke III, S. 418, art. [8.]) aufgezeichneten Reihen
\begin{equation}
\left\{\begin{aligned}
\surd\frac\varpi\pi&=1-2e^{-\pi}+2e^{-4\pi}-2e^{-9\pi}+\ldots,\\[1ex]
\frac12\surd\frac\varpi\pi&=e^{-\frac14\pi}+e^{-\frac94\pi}+e^{-\frac{25}4\pi}+\ldots,
\end{aligned}\right.\tag*{(**)}
\end{equation}
\noindent die den Zusammenhang herstellen zwischen der lemniskatischen Periode und den nach Potenzen von \(e^{-\pi}\) bezw. \(e^{-\frac14\pi}\) fortschreitenden Reihen, deren Exponenten die Quadratzahlen sind. Die auf S. 26 der Scheda stehende Aufzeichnung\begingroup\renewcommand{\thefootnote}{**)}\footnote{Vergl. oben S. 277.}\endgroup
\par\bigskip\noindent X1.\hfill68\hspace*{.14\linewidth}
\par\endgroup

\clearpage\SourceLeaf{567}\ArithmeticHead{538}{L}{TAGEBUCH MIT ERLÄUTERUNGEN.}
\begingroup\small
\indent »Investigatio factorum progressionis infinitae
\[
1+x+x^3+x^6+x^{10}+x^{15}+x^{21}+\text{etc.}=S\text{«}
\]
\noindent zeigt, daß \textsc{Gauss} durch die Entdeckung der beiden Reihen (**) veranlaßt worden ist, seine älteren Untersuchungen über Potenzreihen, deren Exponenten eine arithmetische Progression zweiter Ordnung bilden, zu vergleichen. Von solchen Untersuchungen liegen Spuren vor im art. [9.] der \textit{Exercitat. mathem.} (1796), oben S. 142, in der Tagebuchnotiz Nr. 58 vom 16. Februar 1797, die wieder auf Nr. 7 vom 24. Mai 1796 zurückweist, und in einer mündlichen Überlieferung (Werke III, S. 493, Zeile 3—6), wonach \textsc{Gauss} die Beziehungen zwischen dem arithmetisch-geometrischen Mittel und den Potenzreihen, deren Exponenten Quadratzahlen sind, schon 1794 gekannt haben soll. \textsc{Gauss} konnte also hier einen Zusammenhang zwischen der lemniskatischen Periode und dem arithmetisch-geometrischen Mittel \textls[120]{vermuten}, einen Zusammenhang, der ihm in der Tagebuchnotiz Nr. 98 vom 30. Mai 1799 zur \textls[120]{Gewißheit} geworden ist (\textls[120]{certo} aperietur heißt es dort, gleichsam als Bekräftigung des hier stehenden .\ .\ . nobis aperuit).

Die Reihen (**) hat \textsc{Gauss} in der Tat sofort dazu benutzt, um den Wert von \(\surd\frac\varpi\pi\) auf 26 Dezimalstellen zu berechnen (S. 28 der Scheda, mit Änderung der Reihenfolge abgedruckt Werke III, S. 418, 419, art. [8.]), und auf einem etwa aus derselben Zeit stammenden Zettel (in Fh, Nr. 2, Kapsel 50, demselben, auf dem die Werke III, S. 420, art. [12.] abgedruckten Reihen aufgezeichnet sind) wird diese Zahl ins Quadrat erhoben\footnote{Es heißt dort:
\par\medskip\begingroup\normalsize\centering
Evectio numeri \(0{,}91357913815611682140724\) ad quadratum\par\endgroup
\medskip\noindent und als Ergebnis der Rechnung
\[
\left[\frac\varpi\pi=\right]0{,}8346268416740731872812057352513.
\]
}, also \(\frac\varpi\pi\) vermittelst der ins Quadrat erhobenen Reihen (**) berechnet. Die Vergleichung von
\[
\frac\varpi\pi=4\left(e^{-\frac14\pi}+e^{-\frac94\pi}+e^{-\frac{25}4\pi}+\ldots\right)^2
\]
\noindent mit der z. B. im Leiste (siehe oben S. 176, Gl. [6.]) aufgezeichneten Reihe
\[
4\left(z^{\frac12}+z^{\frac92}+z^{\frac{25}4}+\ldots\right)^2
\]
\noindent liegt dann sehr nahe; sie stimmen für \(z=e^{-\frac12\pi}\) überein.

Als die der Notiz Nr. 95 zugrunde liegenden Entdeckungen können wir also ansehen: Die Aufstellung der trigonometrischen Reihenentwicklungen für den Zähler und Nenner des sinus lemniscaticus, die im wesentlichen mit den \textsc{Jacobi}schen Thetafunktionen für den lemniskatischen Fall übereinstimmen, die Darstellung der lemniskatischen Periode durch Reihen, deren Exponenten die Quadratzahlen sind und möglicherweise hieran anschließend die Wiederaufnahme der Untersuchungen über das arithmetisch-geometrische Mittel (vergl. oben S. 172 ff. und die zugehörigen Bemerkungen S. 260, 261). Auf den von \textsc{Gauss} bei der Entwicklung der Theorie der lemniskatischen Funktionen 1798 eingeschlagenen Weg wirft die Bemerkung ein helles Licht, die \textsc{Gauss} in dem oben S. 248 abgedruckten Briefe an \textsc{Bessel} vom 30. März 1828, in bezug auf \textsc{Abel} macht. Vergl. die Abschnitte III und IV des Aufsatzes »Über \textsc{Gauss}’ Arbeiten zur Funktionentheorie«, Werke X 2.
\par\noindent\hfill\textsc{Klein. Schlesinger.}
\par\endgroup
